\documentclass[11pt]{article}

\usepackage[utf8]{inputenc}
\usepackage[T1]{fontenc}
\usepackage{amsmath, amssymb, amsthm}
\usepackage{geometry}
\geometry{margin=1in}
\usepackage{hyperref}
\hypersetup{
  colorlinks=true,
  linkcolor=blue,
  urlcolor=blue,
}

\title{Toward a Constructive $SU(2)$ Yang--Mills Mass Gap\\
  via a Sheaf-Theoretic Inductive Step}
\author{Arthur Petron}
\date{\today}

\begin{document}
\maketitle

\begin{abstract}
We describe a program for establishing the existence of a mass gap for
$SU(2)$ lattice gauge theory in the continuum limit, organized around an
inductive step on a Wilsonian renormalization-group trajectory. The
inductive state is a sheaf-theoretic package consisting of
(i) a presheaf $\mathcal{A}_k$ of bound-witnesses on a Grothendieck site
over field configurations,
(ii) a polymer-activity sheaf $\mathcal{S}_k$ with Kotecký--Preiss data
$(\rho, K)$, and
(iii) a string-tension lower bound $\widehat{\sigma}_k$ that survives
the block-spin pushforward up to a controlled 2-cell degradation
$\eta_k$ satisfying $\sum_k \eta_k < \widehat{\sigma}_0$. The program
adheres to a strict constructive law: every existence claim returns an
explicit witness with a computable modulus.

This document is the paper skeleton corresponding to Milestone~0.1 of
the accompanying \texttt{work\_plan.md}; the body will be populated as
subsequent milestones close. The authoritative specification lives in
\texttt{docs/outline.md}.
\end{abstract}

\end{document}
